\begin{proofbox}
	\textbf{\textcolor{CProof}{思路提示}}

	利用对称点的两个几何条件：$PP'$ 垂直于直线 $l$，且 $PP'$ 的中点落在 $l$ 上。通过引入比例因子 $k$ 简化计算。

	\medskip

	\textbf{\textcolor{CProof}{正式推导}}

	\begin{description}[style=nextline, leftmargin=2.8em, labelsep=0.5em, itemsep=0.55em]
		\item[\textbf{步骤一（设比例因子）：}]
		      由垂直关系，$PP'$ 的方向与法向量 $(A,B)$ 平行，故存在实数 $k$ 使
		      \[
			      \begin{cases}
				      x' - x_0 = kA, \\
				      y' - y_0 = kB.
			      \end{cases}
		      \]
		      \par\textbf{理由：} 直线 $l$ 的法向量为 $(A,B)$，与 $l$ 垂直的直线方向与法向量相同。

		\item[\textbf{步骤二（中点条件）：}]
		      中点
		      \[
			      M\left(\dfrac{x_0+x'}{2},\dfrac{y_0+y'}{2}\right)
		      \]
		      在 $l$ 上，故
		      \[
			      A\cdot\frac{x_0+x'}{2} + B\cdot\frac{y_0+y'}{2} + C = 0.
		      \]
		      \par\textbf{理由：} 对称点定义要求中点在对称轴上。

		\item[\textbf{步骤三（代入消元）：}]
		      将
		      \[
			      x' = x_0 + kA,\qquad y' = y_0 + kB
		      \]
		      代入中点方程：
		      \[
			      \begin{aligned}
				      A\cdot\frac{x_0 + (x_0+kA)}{2}
				      + B\cdot\frac{y_0 + (y_0+kB)}{2}
				      + C                                               & = 0  \\
				      \Rightarrow\quad
				      Ax_0 + \frac{A^2k}{2} + By_0 + \frac{B^2k}{2} + C & = 0.
			      \end{aligned}
		      \]
		      \par\textbf{理由：} 合并同类项，整理关于 $k$ 的方程。

		\item[\textbf{步骤四（解因子）：}]
		      整理上式为
		      \[
			      (Ax_0+By_0+C) + \frac{A^2+B^2}{2}\,k = 0,
		      \]
		      解得
		      \[
			      k = -\frac{2(Ax_0+By_0+C)}{A^2+B^2}
			      = -2t,
		      \]
		      其中
		      \[
			      t = \dfrac{Ax_0+By_0+C}{A^2+B^2}.
		      \]
		      \par\textbf{理由：} 因为 $A,B$ 不全为零，所以 $A^2+B^2 \neq 0$，方程可解。

		\item[\textbf{步骤五（回代得公式）：}]
		      将 $k = -2t$ 代入步骤一：
		      \[
			      \begin{cases}
				      x' = x_0 + (-2t)A = x_0 - 2At, \\
				      y' = y_0 + (-2t)B = y_0 - 2Bt.
			      \end{cases}
		      \]
		      \par\textbf{理由：} 直接回代即得对称点坐标。
	\end{description}

	\medskip

	\textbf{\textcolor{CProof}{结论回扣}}

	\textbf{如此便得到了对称因子公式。这一推导展示了将几何条件转化为代数比例关系的典型方法。}

	\medskip

	\textbf{\textcolor{CProof}{检验方法}}

	求出 $P'(x',y')$ 后，可以用两个条件快速检验：
	\[
		A\cdot\frac{x_0+x'}{2}+B\cdot\frac{y_0+y'}{2}+C=0,
	\]
	并且
	\[
		(x'-x_0,\ y'-y_0)
	\]
	应与法向量 $(A,B)$ 平行。


	% ---- 仅追加：新增证明方法二、三；原有中点法（方法四）已覆盖，不重复添加 ----
	\medskip
	\textbf{\textcolor{CProof}{方法二（直线方程数值变号法：少计算、易恢复）}}

	\textbf{一句话直觉：} 直线把平面分成两侧。对称前后，两点离直线一样远，却分别在直线两边；所以把点坐标代入同一个直线方程，得到的两个数\textbf{互为相反数}。

	\medskip
	\begin{description}[style=nextline, leftmargin=2.8em, labelsep=0.5em, itemsep=0.6em]
		\item[\textbf{步骤一（把复杂式子记成一个数）：}]
		      记
		      \[
		          f(x,y)=Ax+By+C,\qquad D=f(x_0,y_0)=Ax_0+By_0+C.
		      \]
		      直线 $l$ 上的点满足 $f(x,y)=0$；$D$ 的正、负表示点 $P$ 位于直线的哪一侧。
		      \\[-0.2em]
		      \textbf{注意：}$D$ \textbf{不是}点到直线的实际距离；实际距离为 $|D|/\sqrt{A^2+B^2}$。

		\item[\textbf{步骤二（对称后为什么是 $-D$）：}]
		      设 $P'$ 为对称点、$M$ 为 $PP'$ 的中点。轴对称意味着 $M$ 在 $l$ 上，因此
		      \[
		          f(M)=0.
		      \]
		      由于 $M$ 的两个坐标分别是 $P,P'$ 对应坐标的平均数，代入一次式 $f$ 有
		      \[
		          f(M)=\frac{f(P)+f(P')}{2}.
		      \]
		      因而
		      \[
		          0=\frac{D+f(P')}{2}
		          \quad\Longrightarrow\quad
		          \boxed{f(P')=-D}.
		      \]
		      \textbf{理由：} 对称点到对称轴距离相等且位于相反侧；上面的中点等式也直接验证了数值变号。

		\item[\textbf{步骤三（只能沿垂直方向移动）：}]
		      因为 $PP'$ 垂直于 $l$，而 $(A,B)$ 是直线的法向量，所以只需设一个未知数 $k$：
		      \[
		          \boxed{P'=(x_0+kA,\ y_0+kB)}.
		      \]
		      \textbf{理由：} 沿法向量移动 $k$ 倍，就保证连接两个点的方向垂直于直线；$k$ 可以为正也可以为负。

		\item[\textbf{步骤四（用变号性质求 $k$）：}]
		      把 $P'$ 的坐标代入 $f(P')=-D$：
		      \[
		          \begin{aligned}
		              A(x_0+kA)+B(y_0+kB)+C&=-D,\\
		              \underbrace{Ax_0+By_0+C}_{D}
		                +k(A^2+B^2)&=-D,\\
		              k(A^2+B^2)&=-2D.
		          \end{aligned}
		      \]
		      因 $A^2+B^2>0$，所以
		      \[
		          \boxed{k=-\frac{2D}{A^2+B^2}}.
		      \]
		      \textbf{顿悟：} 为什么有一个 $2$？因为直线方程的值从 $D$ 变为 $-D$，总变化量恰好为 $-2D$。

		\item[\textbf{步骤五（代回对称点坐标）：}]
		      由 $x'=x_0+kA$、$y'=y_0+kB$，得到
		      \[
		          \boxed{\begin{cases}
		              x'=x_0-\dfrac{2A(Ax_0+By_0+C)}{A^2+B^2},\\[7pt]
		              y'=y_0-\dfrac{2B(Ax_0+By_0+C)}{A^2+B^2}.
		          \end{cases}}
		      \]
	\end{description}

	\textbf{初中水平的快速检查：} 对于直线 $y=0$，取点 $P(3,2)$，有 $A=0,B=1,C=0,D=2$。上式给出 $k=-4$，从而
	\[
		P'=(3,2)+(0,-4)=(3,-2),
	\]
	正好是关于 $x$ 轴的对称点。若 $P$ 本来就在直线上，则 $D=0$，公式得到 $P'=P$，也正确。

	\medskip
	\textbf{\textcolor{CProof}{方法三（有向距离法：像照镜子一样理解）}}

	\textbf{一句话直觉：} 要得到镜中的点，先沿垂线走到镜面，再沿同一方向走同样远。总位移就是\textbf{两倍有向距离}，方向由正负号自动决定。

	\medskip
	\begin{description}[style=nextline, leftmargin=2.8em, labelsep=0.5em, itemsep=0.6em]
		\item[\textbf{步骤一（先找到垂直方向）：}]
		      直线 $l:Ax+By+C=0$ 的法向量是 $(A,B)$，它的长度为
		      \[
		          \sqrt{A^2+B^2}.
		      \]
		      将法向量除以长度，就得到长度恰好为 $1$ 的\textbf{单位法向量}：
		      \[
		          \mathbf n=\frac{(A,B)}{\sqrt{A^2+B^2}}.
		      \]
		      \textbf{理由：} $(A,B)$ 只表示方向，长度未必是 $1$；除以自己的长度后，才可以用“走多少个单位长度”来计算位移。

		\item[\textbf{步骤二（把距离写成带正负号的数）：}]
		      点 $P(x_0,y_0)$ 到直线的普通距离是
		      \[
		          \frac{|Ax_0+By_0+C|}{\sqrt{A^2+B^2}}.
		      \]
		      为了同时表示“哪一边”和“离多远”，暂时去掉绝对值，定义\textbf{有向距离}：
		      \[
		          \boxed{d=\frac{Ax_0+By_0+C}{\sqrt{A^2+B^2}}}.
		      \]
		      当 $d>0$ 时，$P$ 位于 $\mathbf n$ 所指的一侧；当 $d<0$ 时，$P$ 位于反方向的一侧。无论哪一侧，实际距离始终是 $|d|$。

		\item[\textbf{步骤三（从原点走到垂足）：}]
		      设 $H$ 是 $P$ 到 $l$ 的垂足。从 $P$ 走到直线，需要抵消有向距离 $d$，所以
		      \[
		          \overrightarrow{PH}=-d\,\mathbf n.
		      \]
		      \textbf{理由：} 若 $d>0$，就逆着 $\mathbf n$ 走 $d$；若 $d<0$，负负得正，就顺着 $\mathbf n$ 走 $|d|$。一个式子处理两侧的点。

		\item[\textbf{步骤四（从垂足再走一次）：}]
		      轴对称的条件是 $H$ 恰好为 $PP'$ 的中点，因此
		      \[
		          \overrightarrow{HP'}=\overrightarrow{PH}=-d\,\mathbf n.
		      \]
		      把两段位移相加：
		      \[
		          \boxed{\overrightarrow{PP'}=-2d\,\mathbf n}.
		      \]
		      \textbf{理由：} 这就是公式中系数 $2$ 的几何来源：从 $P$ 到对称轴走一段，还要从对称轴向另一侧再走同样的一段。

		\item[\textbf{步骤五（两个根号相乘，得到平方和）：}]
		      代入 $d$ 和单位法向量：
		      \[
		          \begin{aligned}
		              \overrightarrow{PP'}
		              &=-2\,\frac{Ax_0+By_0+C}{\sqrt{A^2+B^2}}
		                 \cdot\frac{(A,B)}{\sqrt{A^2+B^2}}\\[5pt]
		              &=-\frac{2(Ax_0+By_0+C)}{A^2+B^2}(A,B).
		          \end{aligned}
		      \]
		      因 $P'=P+\overrightarrow{PP'}$，得到
		      \[
		          \boxed{\begin{cases}
		              x'=x_0-\dfrac{2A(Ax_0+By_0+C)}{A^2+B^2},\\[7pt]
		              y'=y_0-\dfrac{2B(Ax_0+By_0+C)}{A^2+B^2}.
		          \end{cases}}
		      \]
		      \textbf{顿悟：} 分母为什么是 $A^2+B^2$ 而非根号？因为“有向距离”和“单位法向量”分别有一个 $\sqrt{A^2+B^2}$，相乘后根号消失。
	\end{description}

	\textbf{方向再检查：} 仍看直线 $y=0$。对 $P(0,2)$，有 $d=2$，从而向下走 $4$ 得 $P'(0,-2)$；对 $P(0,-2)$，有 $d=-2$，公式会自动改成向上走 $4$ 得 $P'(0,2)$。

	\medskip
	\textbf{\textcolor{CProof}{三种思路如何区分（避免重复背诵）}}

	\begin{itemize}[leftmargin=*, itemsep=0.35em, topsep=0.35em]
		\item \textbf{原证明（中点条件＋法向量）：} 先设未知比例因子，再把中点代入直线方程。这已经覆盖“方法四：中点＋垂直法”，因此不再重复列一遍。
		\item \textbf{方法二（数值变号）：} 先记 $D=Ax_0+By_0+C$，抓住\textbf{$D\to-D$}，一行求出比例因子，适合忘记公式时快速恢复。
		\item \textbf{方法三（有向距离）：} 抓住\textbf{“到镜面一段＋越过镜面一段”}，直接写 $\overrightarrow{PP'}=-2d\mathbf n$，最适合理解几何意义。
	\end{itemize}

	\textbf{最终三行速记：}
	\[
	\begin{gathered}
	D=Ax_0+By_0+C,\qquad S=A^2+B^2,\\
	t=\frac{D}{S},\\
	\boxed{P'=P-2t(A,B)}.
	\end{gathered}
	\]

\end{proofbox}